\section{Saint Constantine the Great}

\begin{quotex}
We are men of the middle ages, not only because that is our destiny, the fatality of history, but also because we will it. You, you are still men of modern times, because you refuse to choose.
\flright{\textsc{Nicolas Berdyaev}}

The greatest miracle of Christianity was succeeding in asserting itself among the European peoples, even taking account the decadence into which numerous traditions of these peoples had plunged. ... Medieval civilisation would be inconceivable, without phenomena such as the appearance of the great Knightly orders, Thomism, a certain mysticism of a high level (e.g., Meister Eckhart), the spirit of the Crusades, etc.
\flright{\textsc{Julius Evola}, \emph{A Justified Pessimism}\url{https://www.gornahoor.net/library/JustifiedPessimism.pdf}}
\end{quotex}

At that time, Rome was ruled by two co-equal Augusti and two subordinate Caesars. That system was unwieldly, as would be expected, since the power hungry competed with each other. A hierarchy must have a single head. That is an analytically true statement and it is modelled on the God-centred cosmic hierarchy.

When Constantine's father --- who was the Augustus of the West --- died, Constantine was declared Augustus by his troops. But he was not the sole claimant. When Maxentius subsequently claimed to be the Augustus, a war with Constantine became inevitable. Eventually, their respective armies were to do battle at Rome although Constantine was greatly outnumbered. Before the battle, there appeared in the sky a large cross of light above the Sun bearing the inscription, ``\emph{In Hoc Signo Vinces}'' (``with this sign, you will conquer''). This was a public event, witnessed not only by Constantine, but also by his army.

This public sign was followed by a private apparition of Lord Jesus to Constantine in a dream. Christ explained to him how to prevail in the battle. He also ordered Constantine to construct a \emph{laborum} with the letters Chi Rho affixed. This sign was also placed on the shields and arms of the soldiers.

As prophesied, Maxentius was defeated. Constantine issued the Edict of Milan, which did not make Christianity the state religion, but rather established freedom of religion; this ended three centuries of Christian persecution. Pagan gods would be worshiped at the Triumphal Arch built to commemorate Constantine's victory. In this way, Christianity eventually prevailed in the Empire due to its superiority over paganism, which had fallen into decadence. Constantine himself was not baptized until several years later.

So you are left with a choice. You can pretend none of this happened and that there was nothing but natural events. That is the position of scholars and most moderns. However, that does not explain the
\emph{laborum}, which has not secular source; it was not a Roman or pagan symbol.

Or you can accept that Providence works in human events and that Christ intervened in to assure Constantine's victory. There is the greater holy war within oneself; but there is also the lesser holy war, which is also fought on the material plane in the name of Christ. The Christian state and the holy war are legitimate manifestations in Christendom. Any denial of this is actually some other religion. Hence, there is no Christian pacifism. The existence of Knightly orders, the Templars, Crusades attest to that. It makes no sense to assert that the first 12 centuries of Christianity was not Christian. If doctrine is ``up for grabs'', i.e., not more than personal opinion, it would hardly be worth following.

Sophisticates have tried to impugn Constantine's character, some 17 centuries after the fact. But consider how varied are the judgments of contemporary public figures and you will never again take such speculations across the centuries seriously. Moreover, the mindset of that time is opaque to most of our contemporaries.

Speaking what is good is better than silence, and silence is better than speaking evil.


\flrightit{Posted on 2019-12-16 by Cologero}

\begin{center}* * *\end{center}

\begin{footnotesize}
\begin{sffamily}
\texttt{Mate on 2019-12-17 at 02:04 said:}

I fail to see how the described miracle of Constantine and the christian turn that it moved him towards connects to the Crusades being legitimate `holy wars'. Ironically if we look at the 4th one it sacked the city of Constantine and made a bloodbath among its people.

\hfill

\texttt{Dennis on 2019-12-17 at 10:40 said:}

Mate: One thing people must keep in mind regarding the Crusades (despite occasional unfortunate episodes within them) is that they originated as defensive actions to defend Christians after centuries on persistent aggression from Muslims aggression and Muslim imperialism. It is only in the past 100 years or so that self-hating Westerners have imbibed the anti-Christian, anti-Crusade notion that they were simply a manifestations of Christian, European ``imperialism'', ``colonialism,'' etc. (insert the evil du jour the PC crowd denouncing). The First Crusade especially was an entirely justified Christian Holy War. Thomas Madden is perhaps the best modern historian of the Crusades who has set the record straight.

\hfill

\texttt{Boreas on 2019-12-17 at 10:59 said:}

Mate, I think this refers to the Spirit of Crusades being sacred, not necessarily everything that happened on a material plane.

``History is meaningful only to the degree it follows Myth.''

\hfill

\texttt{Cologero on 2020-01-02 at 09:09 said:}

History, as officially taught, limits itself to exterior events, which are only the effects of something deeper; and it sets these events forth in a tendentious manner under the influence of all the modern prejudices. And further, there is a veritable monopoly on historical studies in the interest of parties, both political and religious.
\flright{\textsc{Rene Guenon}}

See \textit{The Truth is far from Suspected}\footnote{\url{https://www.gornahoor.net/?p=5559}}.

If one does not believe in chance, one is forced to admit the existence of some kind of equivalent of an established plan, but one which evidently does not need to be formulated in any document.

\hfill

\end{sffamily}
\end{footnotesize}

